\section{Equivariant spatial observables}\label{sec:observables}
The shell selects a polar vertical \(e_z\), outward normals \(n_i\), tangents \(t_i=e_z\times n_i\) and centre o. For the triad write
\begin{gather*}
\Omega=q+ip,\quad e=(1,1,1)^T,\quad
q=ae+u,\ p=be+v,\quad u,v\perp e,\\
N_f=[n_A\ n_B\ n_C],\qquad T_f=[t_A\ t_B\ t_C].
\end{gather*}
Exact coordinates appear in Appendix~\ref{app:notation}. The frame matrices have Gram matrix \((3/2)\Pi\), \(\Pi=I-ee^T/3\). \(N_f\) is not ring size.

Let P cycle A,B,C and \(Q_v\) swap A,C. Polar covariance of the adopted decoder gives the actual state actions
\begin{equation}
C_3:(q,p)\mapsto(Pq,Pp),\quad
\sigma_v:(q,p)\mapsto(-Q_vq,Q_vp),\quad
\sigma_h:(q,p)\mapsto(q,-p).\label{eq:reps}
\end{equation}
They are not ordinary permutation of both channel triples. Polar free vectors transform by G, axial vectors by \(\det(G)G\), and points affinely about o.
This is the kinematic representation at every state. Dynamical covariance carries the coefficient triple k with the face permutation; at fixed unequal k only its stabilizer remains. Mirrors involving common phase require the regular pre-synchronization domain for nonzero lambda. The census does not assume a global dynamical symmetry across the zero stratum.

\begin{theorem}[Complete ambient census]\label{thm:observables}
For real polynomial state maps equivariant under full \(D_{3h}\), the linear polar basis is \(T_fu,be_z\), and the homogeneous quadratic polar basis is
\begin{equation}
aN_fu,\quad N_f(u\odot u),\quad bN_fv,\quad
N_f(v\odot v),\quad\Gamma_{\rm ch}e_z,\qquad
\Gamma_{\rm ch}=e^T(u\times v).\label{eq:polarbasis}
\end{equation}
The axial linear basis is \(T_fv,ae_z\), and the quadratic axial basis is
\(aN_fv,bN_fu,N_f(u\odot v)\). No nonzero constant ambient equivariant exists.
\end{theorem}
\begin{proof}
The state representation on a,u,b,v decomposes as \(A_2'\oplus E'\oplus A_2''\oplus E''\). The polar target is \(E'\oplus A_2''\), the axial target \(E''\oplus A_2'\). Direct expansion of q-q, p-p and q-p products gives
\begin{equation}
\operatorname{Sym}^2\mathcal X_\R=
4A_1'\oplus4E'\oplus2A_1''\oplus A_2''\oplus3E''.\label{eq:sym2}
\end{equation}
Each unmixed sector gives two invariant scalars and two transverse copies; the mixed sector gives two mirror-odd scalars, the oriented scalar \(\Gamma_{\rm ch}\), and three transverse copies. Counting matching irreducibles gives dimensions 2,5 and 2,3. The listed maps obey the generator laws. Independent means and distinct q-q, p-p and q-p monomials prove independence, hence completeness. Neither target contains the trivial representation, so constant maps vanish.
\end{proof}

An additional demand of common-phase invariance still leaves the three-dimensional quadratic polar space
\begin{equation}
\operatorname{span}\{N_f(au+bv),\,N_f(u\odot u+v\odot v),\,
\Gamma_{\rm ch}e_z\}.\label{eq:U1}
\end{equation}
The axial invariant space is \(\operatorname{span}\{W\}\), \(W=T_f(q\times p)\), and no invariant linear map survives. Apply the common-phase generator
\(\sum_i(-p_i\partial_{q_i}+q_i\partial_{p_i})\) to the complete bases: its kernel is exactly the displayed dot-product and oriented-area combinations. This is an optional kinematic condition, not a global phase quotient of F at zero.

\subsection{Chiral projections and native observable changes}
\begin{equation}
C_{\rm ch}=q\times p=e\times(av-bu)+u\times v,\quad
W=\sqrt3(bN_fu-aN_fv),\quad \Gamma_{\rm ch}=e^TC_{\rm ch}.\label{eq:chiral}
\end{equation}
Raw C transforms as \(P,Q_v,-I\), neither ambient polar nor axial law. Its transverse part gives horizontal axial W and its common part gives vertical polar \(\Gamma_{\rm ch}e_z\). Each is unique up to scale within its corresponding antisymmetric sector, not among all observables. Together
\[
C_{\rm ch}=\tfrac23T_f^TW+\tfrac13\Gamma_{\rm ch}e.
\]
For pure transverse independent u,v, W vanishes but Gamma does not. Axial parity alone need not mean chirality: \(ae_z\) is nonzero on common purely real states.

Equal channels kill every quadratic spatial map. Unequal amplitudes at common phase can yield intensity/coherence vectors while \(C_{\rm ch}=0\). At \(\Omega=r_0e+iv\), \(v\perp e\), the first-order axial readout is \(-\sqrt3r_0N_fv\). For \(u=(1,-1,0),v=(1,1,-2),a=b=0\), one has \(C_{\rm ch}=2e,\Gamma_{\rm ch}=6,W=0\).

An observable V receives no additional evolution:
\[
\Delta V=V(F(\Omega))-V(\Omega).
\]
For instance \(\Delta(T_fu)=T_f\Delta u\) and
\[
\Delta\Gamma_{\rm ch}
=e^T(u\times\Delta v+\Delta u\times v+\Delta u\times\Delta v).
\]
All quadratic increments follow substitution; local response is \(DV(\Omega_*)DF(\Omega_*)\). At a positive equal triad, radial mean, radial transverse, imaginary mean and imaginary transverse factors are
\(1-2\eps k,\ 1-2\eps k-3g,\ 1,\ (1-3g)(1-9\lambda)\).
On unequal branches use \eqref{eq:SP}. Existing one-step examples generate basis observables from zero (Appendix~\ref{app:closure}); nonzero readouts at fixed states further show why a vector value is not displacement per step.

Symmetry therefore does not select a unique state-dependent core position or motion observable. A physical position law remains an additional ingredient. \src{CM0 \S\S2--6, L28--L33.}
\begin{figure}[htbp]\centering
\includegraphics[width=\linewidth]{../figures/04_observables.pdf}
\caption{Independent chiral projections and the larger equivariant spaces. Axial denotes reflection parity, not vertical direction. Multiplicities count polynomial-map spaces, not physical modes. Source: CM0 \S\S3--4 and \eqref{eq:polarbasis}--\eqref{eq:chiral}.}
\end{figure}

\section{State--shell attachment and finite-carrier multiplicity}\label{sec:attachment}
The accepted fixed mesh has three octagonal faces, three vertical seams and two nine-edge rim cycles. A face lies in the affine plane \(c_i+W_i\), \(W_i=n_i^\perp\); \(W_i\) is unbounded.

\begin{definition}[Adopted decoder]
\begin{equation}
D_i(q_i+ip_i)=q_it_i+p_ie_z,\qquad
E_i(v_i)=t_i\cdot v_i+i e_z\cdot v_i,\qquad
D:\C^3\to\bigoplus_iW_i.\label{eq:decoder}
\end{equation}
\end{definition}
Orthonormality proves \(ED=I\), \(DE=\diag(I-n_in_i^T)\) on ambient arrays, and \(\sum\|D_i\Omega_i\|^2=\sum|\Omega_i|^2\). Basing at the adopted centre gives \((c_i,v_i)\); the endpoint \(c_i+v_i\) need not lie in the finite face and is not a particle-position law.

\begin{proposition}[Discrete connection and conjugacy]
The maps \(T_{ij}=t_it_j^T+e_ze_z^T:W_j\to W_i\) are orthogonal isomorphisms satisfying \(T_{ij}T_{jk}=T_{ik}\). They define a flat discrete rank-two connection on the face-label graph. The exact conjugate update is
\begin{align}
\widetilde v_i&=v_i+\eps(k_i-\|v_i\|^2)v_i
 +g\sum_{j\ne i}(T_{ij}v_j-v_i),\\
v_i^+&=\cos\delta_i\,\widetilde v_i+
 \sin\delta_i\,n_i\times\widetilde v_i,\qquad F_{\rm face}=DFE.\label{eq:faceupdate}
\end{align}
\end{proposition}
\begin{proof}
T carries the two frame coefficients unchanged, proving composition and isometry. Closed transport is identity on its tangent domain, a projector ambiently. Encoding the first line recovers the native pre-stage. The identity \(E_i(n_i\times v)=iE_i(v)\) turns the second into phase multiplication, including the zero value. D/E invert on their domain.
\end{proof}
The implementation steps stored canonical Omega and then decodes; it does not repeatedly encode rounded displays. No second recurrence or carrier motion is added.

A section on three labelled fibres differs from a field on each continuous polygon, seam/rim traces, shell deformation and physical position. Only the first is natively defined. Conditional constant face extensions illustrate the distinction: ambient equality at every seam forces \(q=0,p=be\), because distinct tangent planes intersect only in \(\R e_z\). Matched parallelism instead forces \(q=ae,p=be\). Neither criterion is imposed on arbitrary native states.

\subsection{Multiplicity on finite carriers}
\begin{center}\small
\begin{tabular}{lrrr}\toprule
Carrier & degree 0 & degree 1 & degree 2\\\midrule
Ordinary face scalar or normal coefficient &1&1&8\\
Polar tangent vector per face &0&4&8\\
Polar ambient vector per face &1&5&16\\
Polar centre vector &0&2&5\\
Axial centre vector &0&2&3\\
Ordinary scalar per unordered seam &1&1&8\\
Polar vertical tangent vector per seam &0&2&4\\
Polar ambient vector per unordered seam &1&5&16\\
Scalar per whole rim &1&1&5\\
Polar ambient vector per whole rim &1&3&12\\\bottomrule
\end{tabular}
\end{center}
The proof uses \eqref{eq:sym2} and the target representations; complete independent bases appear in Appendix~\ref{app:census}. Each seam is indexed equivariantly by its opposite face. Horizontal reflection exchanges the two rims. No individual seam or boundary point is selected.

Already linear tangent assignments have four independent constants:
\[
Q=\alpha ae+\beta u,\qquad P=\gamma be+\delta v,\qquad
v_i^{\rm cand}=Q_it_i+P_ie_z.
\]
D adopts all four as one. A normal coefficient \(f_i n_i\) does not specify compatible vertex displacement; a whole-loop vector need not be tangent to every segment. Pointwise fields need a spatial function space beyond this census.

\subsection{Preparation and boundary types}
The exact type chain is
\begin{equation}
(R,d;\xi\in\C^2)\mapsto\theta,\psi_0\in\C^2\otimes\C^3
\xrightarrow{U^n}\psi_n\xrightarrow{\chi^\dagger\otimes I}\Omega_0
\xrightarrow{D}\bigoplus_iW_i.\label{eq:prep}
\end{equation}
The named toy response uses
\[
\theta=\arccos(d/2R),\quad
G(\theta)=\sqrt{(2\theta-\sin2\theta)/\pi},\quad
\psi_0=G(\theta)(\xi\otimes f_0),\quad f_0=e/\sqrt3,
\]
with \(R>0,0\le d\le2R\). The fixed transfer is \(U=B_s\otimes A_s\), with named normalized eigenbra \(\chi^\dagger B_s=b_\chi\chi^\dagger\). Its established reduction is
\[
\Omega_0=G(\theta)(\chi^\dagger\xi)b_\chi^n h_n f_j,\quad
h_n=(r_sq_s^2c_s^2)^m(1,q_s,q_s^2c_s)_j,\quad n=3m+j,
\]
where \(q_s=e^{-.423},r_s=e^{.577},c_s=1-.618\) are fixed initialization literals; \(f_j\) are the ordered Fourier vectors.

The final Omega can use D. No map attaches preparation internals to material rims, seams, faces or scaffold cells. The complex six-state space has twelve real dimensions; the incident pair is not two rim traces; the lens chart is not computed from shell geometry. Source signatures and imports use them for initialization only; F receives Omega and its native parameters. The reference scaffold is a separate fixed family of closed sets with a special explicit Paper-C comparison, not a state attachment to gaps. There is no native boundary/seam feedback. \src{SA0 \S\S2--7; prior Paper B/Atlas 04 and Atlas 08, L34--L37.}
\begin{figure}[htbp]\centering
\includegraphics[width=\linewidth]{../figures/05_typed_attachment.pdf}
\caption{Typed attachment. D is the existing invertible representation; preparation reaches it through final Omega. Dashed absent arrows mark missing definitions, not proposed interactions. Source: SA0 \S\S2--3,7 and \eqref{eq:prep}.}
\end{figure}

\section{Closure hierarchy and coherence quotient}\label{sec:closure}
\begin{theorem}[Fibre closure criterion]\label{thm:closure}
For maps \(\mathcal A:X\to Y\), \(F:X\to X\), an induced map satisfying \(\mathcal AF=F_{\mathcal A}\mathcal A\) on \(\mathcal A(X)\) exists iff
\begin{equation}
\mathcal A(x)=\mathcal A(x')\ \Longrightarrow\
\mathcal A(Fx)=\mathcal A(Fx').\label{eq:fibre}
\end{equation}
\end{theorem}
\begin{proof}
Necessity is composition with \(F_{\mathcal A}\). For sufficiency define \(F_{\mathcal A}(y)=\mathcal A(Fx)\) using any representative; the implication makes it well defined.
\end{proof}
This asserts no locality or smoothness. A restricted-domain claim requires the same implication on every fibre in that domain.

The complete decoded/based tangent triple closes globally by bijectivity of D. One face generally does not: with \(\eps=1/20,g=1/5,k=e\), real inputs \((1,1,1)\) and \((1,2,1)\) have equal A coefficient but next coefficients 1 and \(6/5\). Pre-states are positive, so this holds for all lambda.

\subsection{Lower-readout obstructions}
At the same eps,g,k, \((1,1,1)\) and \((1,-1,1)\) have identical intensities but next intensities
\[
(1,1,1)\quad\hbox{and}\quad(9/25,1/25,9/25).
\]
Synchronization preserves moduli; total intensity also differs. At g=0 the explicit exception is \(I_i^+=I_i[1+\eps(k_i-I_i)]^2\).

For chirality, set lambda=0 and compare
\[
(q,p)=(e,(1,-1,0)),\qquad(q',p')=(2e,(1,-1,0)/2).
\]
Both give initial \(C_{\rm ch}=(1,1,-2)\), but next values
\[
(7/20,7/20,-133/200),\qquad
(323/1600,323/1600,-1273/3200).
\]
Their common and transverse components differ. C, W, \(\Gamma_{\rm ch}\) and their joint chiral representation therefore fail general closure. These exact fibre counterexamples do not exclude special parameter or invariant-subset exceptions.

\subsection{Complex coherence}
Distinguish the response matrix \(H_r\) from the coherence
\[
\mathcal H=\Omega\Omega^\dagger=S_G-iA_G,\qquad
S_G=qq^T+pp^T,\quad A_G=qp^T-pq^T .
\]
For rank-at-most-one positive semidefinite Hermitian \(\mathcal H\), put
\begin{equation}
M(\mathcal H)=I+\eps\diag(k-\diag\mathcal H)+g\Delta_3,\qquad
\widetilde{\mathcal H}=M\mathcal HM^T.\label{eq:Hpre}
\end{equation}
\begin{theorem}[Regular complex-coherence closure]\label{thm:Hclosure}
Where all \(\widetilde{\mathcal H}_{ii}>0\), set
\[
\rho_i=\sqrt{\widetilde{\mathcal H}_{ii}},\quad
\delta_i=\lambda\sum_{j\ne i}\im
\left(\frac{\widetilde{\mathcal H}_{ji}}{\rho_j\rho_i}\right)^3,\quad
U_\delta=\diag(e^{i\delta_i}).
\]
The exact induced update is
\begin{equation}
\boxed{\mathcal H^+=U_\delta\widetilde{\mathcal H}U_\delta^\dagger.}\label{eq:Hpost}
\end{equation}
\end{theorem}
\begin{proof}
\(M\Omega\) is the native pre-state. Its normalized ji coherence entry is \(e^{i(\widetilde\phi_j-\widetilde\phi_i)}\). Cubing supplies the native sine increment. The final phase multiplication gives \eqref{eq:Hpost}; everything depends only on \(\mathcal H\). Equivalently its fibres are common-phase orbits, and F is common-phase equivariant on this domain.
\end{proof}
Iteration is justified while the regularity condition persists; the domain is not claimed invariant.

\begin{theorem}[Regular real-Gram closure]\label{thm:Sclosure}
The real Gram matrix \(S_G\succeq0\), of rank at most two, also closes where \((MS_GM^T)_{ii}>0\).
\end{theorem}
Here \(M=I+\eps\diag(k-\diag S_G)+g\Delta_3\); its dependence is only on the retained diagonal.
\begin{proof}
Equal real Gram factors \(QQ^T=Q'Q'^T\) have identical pairwise row inner products. The row-span identification is an isometry and extends to an \(O(2)\) transformation, giving \(Q'=QO\). Such changes are common phase and/or conjugation of Omega. The real-parameter F is conjugation equivariant, and on this domain common-phase equivariant. Output real Grams therefore agree, proving the fibre criterion, including rank-deficient factors.
\end{proof}
S loses signed orientation; closure survives because the discarded distinction is a symmetry. This is a nonlocal channel quotient, not one scalar per face.

\subsection{The essential zero-stratum exception}
At \(\eps=g=0,\lambda=1/10\), let \(z=(1,-1,0)\), \(z'=e^{i\pi/6}z\). They have identical complex and real Grams. For z all sine increments vanish; for z' the two nonzero increments are \(-\lambda,+\lambda\), since the zero phase remains zero. Thus
\begin{equation}
\mathcal H^+_{12}=-1,\quad(\mathcal H')^+_{12}=-e^{-i/5},\qquad
S^+_{G,12}=-1,\quad(S'_G)^+_{12}=-\cos(1/5).\label{eq:zerocounter}
\end{equation}
They differ since \(0<1/5<\pi\). Neither quotient closes globally for the native nonzero-lambda convention. Both congruence updates are global when lambda=0. No repaired convention is proposed.

The hierarchy is exact: D retains full state; lower face/intensity/area readouts generally discard necessary information; coherence admits a regular symmetry quotient; the native zero convention prevents unrestricted extension. Fixed geometry is trivially constant; preparation receipts, clocks and observer memory are other data types, not autonomous Omega-only shell variables. \src{SA0 \S8, L38--L42.}
